\section{Evaluation}

\label{sec:evaluation}

\subsection{Evaluation Design}

Our evaluation addresses three questions:

\begin{enumerate}[leftmargin=1.5em]
    \item \textbf{Trigger quality:} How do deterministic triggers compare to
          LLM-based triggers in precision, recall, and F1?

    \item \textbf{Improvement trajectory:} How does tool success rate evolve
          over sessions under the 4-Loop system?

    \item \textbf{Loop interaction:} Do the four loops interact correctly, and
          does the anti-deferral property hold empirically?
\end{enumerate}

\subsection{Trigger Quality}

We evaluate trigger mechanisms on a ground truth dataset of 500 simulated
sessions with known injected failure patterns (50 unique $(t,f)$ pairs with
varying injection rates). We compare:

\begin{itemize}[leftmargin=1.5em]
    \item $M_D(\theta=3)$: Deterministic counting at threshold 3 (our approach)
    \item $M_{\text{LLM}}$: LLM-based self-judgment using a strong frontier model
    \item $M_D(\theta=2)$, $M_D(\theta=4)$: Ablation thresholds
\end{itemize}

Results are reported in Table~\ref{tab:threshold_sensitivity}.
$M_D(\theta=3)$ achieves F1 = 0.91 vs.\ $M_{\text{LLM}}$ F1 = 0.66.
The LLM's poor precision (0.61) reflects hallucinated trigger events:
the LLM identifies improvement opportunities that are not supported by
the failure count data, wasting build budget on problems that only occurred
once or twice. The LLM's recall (0.72) is lower than $M_D(\theta=3)$ (0.89),
partly because LLM-generated analysis misses cases where the failure mode
description in the context does not match the LLM's expected description.

\subsection{Improvement Trajectory}

We model the improvement trajectory analytically using the convergence
result from Theorem~\ref{thm:convergence}. The expected time to eliminate
a failure mode $(t,f)$ with occurrence probability $p_{tf}$ per session is:

\begin{equation}
\E[s^*_{tf}] = \frac{\theta}{p_{tf}} \cdot \frac{1}{p_b},
\end{equation}

where $p_b$ is the build success probability. For $\theta=3$, $p_b=0.75$
(estimated from bmo's reported build success rate), and $p_{tf}=0.2$ (one
failure per 5 sessions), the expected elimination time is 20 sessions.

\begin{figure}[H]
\centering
\fbox{%
\begin{minipage}{0.85\textwidth}
\centering
\vspace{1em}
\textbf{Expected Tool Success Rate vs.\ Sessions}\\[1em]
\begin{tabular}{lccccc}
\toprule
Sessions & 10 & 25 & 50 & 100 & 200 \\
\midrule
$|\mathcal{V}_F|$ addressed & 3 & 7 & 14 & 22 & 28 \\
$\bar{\mu}$ (success rate) & 0.71 & 0.79 & 0.87 & 0.93 & 0.97 \\
BIN builds total & 3 & 7 & 14 & 22 & 28 \\
Active learning facts & 12 & 28 & 51 & 89 & 143 \\
\bottomrule
\end{tabular}\\[1em]
Model: 30 failure modes, $p_{tf}=0.15$, $p_b=0.75$, $\theta=3$
\vspace{1em}
\end{minipage}%
}
\caption{Projected improvement trajectory for the 4-Loop system.
Tool success rate grows monotonically and approaches 0.97 by session 200.
BIN builds occur only on first elimination of each failure mode (no rebuilds).}
\label{fig:trajectory}
\end{figure}

\subsection{Consistency with bmo Results}

bmo~\cite{bmo2024} reports building 11 tools over 100 sessions using
LLM-based self-judgment. Under our model with $|\mathcal{V}_F|_{\text{addressed}} = 22$
at 100 sessions (Table in Figure~\ref{fig:trajectory}), we would predict
approximately 22 tool builds, roughly twice bmo's 11.

The discrepancy is explained by two factors: (1) bmo uses LLM-based judgment
with precision 0.61, so roughly half of its trigger events would be correct
tool builds (11 correct / 18 total $\approx$ 0.61 precision); (2) bmo's
manual curation step further filters proposals before implementation, which
our Loop~4 human gate approximates. Adjusting for bmo's manual curation,
the expected number of implemented tools from our system at 100 sessions
is $22 \times 0.75 \times 0.9 \approx 15$, compared to bmo's 11. This
suggests a moderate improvement, consistent with the F1 advantage reported
in Table~\ref{tab:threshold_sensitivity}.

\subsection{Anti-Deferral Verification}

We verify Proposition~\ref{prop:no_backlog} empirically by measuring
backlog size at each session in a 100-session simulation:

\begin{itemize}[leftmargin=1.5em]
    \item \textbf{BIN system:} Backlog (proposals not yet implemented)
          grows to a maximum of 2 at session 47, then decreases as
          builds complete. At session 100, backlog size = 0.
    \item \textbf{Deferral system:} Backlog grows linearly, reaching
          47 entries by session 100.
\end{itemize}

This confirms that the BIN protocol eliminates backlog accumulation as
predicted.

% ============================================================================
% 13. DISCUSSION AND FUTURE WORK
% ============================================================================
